% 张三慧《大学物理学》第4版，第二卷，第19章《磁场中的磁介质》
% 主教材扫描 PDF 物理页187--189，章末习题19.1--19.14。
% 图19.25--19.27与题目同属源材料；此处为学术审核前的重述。
\section*{第19章\quad 磁场中的磁介质}
\subsection*{习题19.1—19.14}

\paragraph{习题 19.1}
顺磁样品每单位体积 \(N\) 个原子，各有固有磁矩大小 \(m\)，在外场
\(B\) 中仅允许磁矩平行、反平行磁场，势能分别为 \(-mB\)、\(+mB\)。
用玻尔兹曼权重证明磁化强度
\[
 M=Nm\frac{e^{mB/(k_BT)}-e^{-mB/(k_BT)}}
 {e^{mB/(k_BT)}+e^{-mB/(k_BT)}}.
\]
推得高温弱场 \(M=Nm^2B/(k_BT)\)（居里定律），
以及低温强场饱和 \(M=Nm\)。

\paragraph{习题 19.2}
在图19.2所述电子轨道角动量 \(\boldsymbol L\)
与外场 \(\boldsymbol B\) 夹角 \(\theta\) 时：
（1）证明磁场对电子轨道的力矩大小为
\(BeL\sin\theta/(2m_e)\)；
（2）证明轨道角动量在场中进动角速度
\(\omega_p=eB/(2m_e)\)，求 \(B=1\,\mathrm T\) 的进动角速度。

\paragraph{习题 19.3}
在经典玻尔氢原子模型里，基态电子轨道半径
\(r=0.53\times10^{-10}\,\mathrm m\)，速率
\(v=2.2\times10^6\,\mathrm{m/s}\)。
（1）求轨道电流在圆心的磁场；
（2）电子自旋角动量
\(S=\hbar/2=0.53\times10^{-34}\,\mathrm{J\,s}\)，
自旋磁矩大小
\(1.41\times10^{-26}\,\mathrm{A\,m^2}\)，
求它绕轨道磁场的进动角速度。
\textit{按主教材给定的经典轨道/自旋模型计算。}

\paragraph{习题 19.4}
假设铁晶体每个铁原子均有两个电子自旋参与磁化且全部平行。
铁棒为直径 \(1.0\,\mathrm{cm}\)、长 \(12\,\mathrm{cm}\) 的圆柱，
密度 \(7.8\,\mathrm{g/cm^3}\)，摩尔质量
\(55.85\,\mathrm{g/mol}\)。
计算：（1）参与排列的电子数；（2）总磁矩；
（3）绕该棒轴线均匀的等效面电流大小；
（4）其内部由该面电流产生的磁场大小。
\textit{每个电子自旋磁矩取玻尔磁子 \(\mu_B\)。}

\paragraph{习题 19.5}
以长 \(8.5\,\mathrm{cm}\)、宽 \(1.0\,\mathrm{cm}\)、
厚 \(0.02\,\mathrm{cm}\) 的铁制磁针为模型，
所有参与磁化的自旋全部同向；铁密度
\(7.8\,\mathrm{g/cm^3}\)，摩尔质量
\(55.85\,\mathrm{g/mol}\)，每个铁原子按两电子自旋。
求磁矩、在地磁场
\(B_E=0.52\times10^{-4}\,\mathrm T\)
中磁矩与场垂直时的力矩，以及磁矩与该场平行时的势能。

\paragraph{习题 19.6}
螺线环的中心周长 \(l=10\,\mathrm{cm}\)，
线圈 \(N=20\) 匝，电流 \(I=0.1\,\mathrm A\)。
（1）计算无介质时环内 \(B_0,H_0\)；
（2）环内填满相对磁导率 \(\mu_r=4200\) 的材料，
计算 \(B,H\)；（3）分离自由电流贡献 \(B_0\)
和磁化面电流贡献 \(B'\)。

\paragraph{习题 19.7}
铁螺线环平均周长 \(30\,\mathrm{cm}\)，
截面积 \(1.0\,\mathrm{cm^2}\)，绕 \(300\) 匝，
通电 \(0.032\,\mathrm A\)。
环内磁通为 \(2.0\times10^{-6}\,\mathrm{Wb}\)。
计算 \(B,H\)、磁化面电流密度、
磁导率、相对磁导率与磁化强度 \(M\)。

\paragraph{习题 19.8}
铁磁化特性实验使用环形螺线管，1000匝，
平均半径 \(15.0\,\mathrm{cm}\)，电流 \(2.0\,\mathrm A\)，
测得环内磁感应强度 \(B=1.0\,\mathrm T\)。
求环内场强 \(H\)、材料磁导率和相对磁导率、
磁化面电流密度。

\paragraph{习题 19.9}
教材图19.25为铁芯初始磁化 \(B\)-\(H\) 曲线。
长直螺线管带铁芯时通 \(I=6.0\,\mathrm A\)
可在管内得到 \(B=1.2\,\mathrm T\)。
抽走铁芯但仍要相同 \(B\)，导线电流应增大到多少？
\textit{须根据源图19.25读取原来 \(H\) 后计算。}

\paragraph{习题 19.10}
参照图19.25初始磁化曲线。
长直螺线管每米800匝，
要使环内铁芯 \(B=1.0\,\mathrm T\)，
估算所需电流。
\textit{图像插值的读取精度须注明。}

\paragraph{习题 19.11}
图19.26为矫顽力 \(H_c=2\,\mathrm{A/m}\)
的矩形磁滞回线材料制作的环形磁芯，
外直径 \(0.8\,\mathrm{mm}\)，内直径
\(0.5\,\mathrm{mm}\)，高度 \(0.3\,\mathrm{mm}\)。
若原来磁化方向如图，欲使磁化方向完全翻转，
求脉冲电流最小峰值；按磁路长度与矫顽力关系计算。
\textit{原图19.26中的磁化方向与通电方向需保留审核。}

\paragraph{习题 19.12}
铁环平均周长 \(61\,\mathrm{cm}\)，
其中一段长 \(1\,\mathrm{cm}\) 为空气隙，
绕线 \(N=1000\) 匝。
当电流 \(1.5\,\mathrm A\) 时，
气隙中测得 \(B=0.18\,\mathrm T\)。
忽略漏磁与气隙边缘磁场发散，求铁芯的相对磁导率。

\paragraph{习题 19.13}
电磁铁铁芯几何如教材图19.27，
铁芯中线长度 \(l_1=500\,\mathrm{mm}\)，
空气隙总长 \(l_2=20\,\mathrm{mm}\)，
铁芯相对磁导率 \(\mu_r=5000\)。
空气隙中期望磁场 \(B=3\,\mathrm T\)，
求线圈安匝数 \(NI\)。
\textit{气隙中磁通不发散，铁芯和气隙近似同 \(B\)。}

\paragraph{习题 19.14}
电铃的线圈和铁芯：
铁芯磁路 \(14.4\,\mathrm{cm}\)，
空气隙宽 \(2.0\,\mathrm{mm}\)，
截面积 \(0.60\,\mathrm{cm^2}\)，
铁芯相对磁导率 \(\mu_r=1600\)。
要求通过空气隙的磁通量 \(4.8\times10^{-6}\,\mathrm{Wb}\)，
求安匝数 \(NI\)。
若线圈两端电压 \(220\,\mathrm V\)，
功率 \(20\,\mathrm W\)，
求线圈匝数 \(N\)。
\textit{忽略磁通泄漏及铁芯的磁滞与饱和。}
